% ===================================================================
% PREAMBOLO DEGLI APPUNTI
% Questo file contiene tutti i pacchetti e le impostazioni
% principali del documento.
% ===================================================================

% -------------------------------------------------------------------
% LINGUA E CODIFICA
% -------------------------------------------------------------------
\usepackage[italian]{babel}      % Gestione della lingua italiana
\usepackage[T1]{fontenc}         % Codifica dei font (permette di usare le accentate)
\usepackage[utf8]{inputenc}      % Codifica del file di input

% -------------------------------------------------------------------
% TIPOGRAFIA E MICRO-TIPOGRAFIA
% -------------------------------------------------------------------
\usepackage{microtype}           % Migliora la spaziatura del testo
\usepackage{setspace}            % Per gestire l interlinea
\onehalfspacing                  % Interlinea 1.5

% Spaziatura tra i paragrafi
\AtBeginDocument{
  \setlength{\parskip}{0.35\baselineskip}%
}

% -------------------------------------------------------------------
% LAYOUT DELLA PAGINA (KOMA-Script)
% -------------------------------------------------------------------
\usepackage[automark,headsepline]{scrlayer-scrpage} % Intestazioni e pi di pagina avanzati
\KOMAoptions{headinclude=false,footinclude=false}    % Bilanciamento margini KOMA
\clearpairofpagestyles
\ihead{\headmark}  % Titolo del capitolo a sinistra
\ohead{\pagemark}   % Numero di pagina a destra
\pagestyle{scrheadings}

% Margini
\usepackage[
    a4paper,
    top=2.5cm,    % Margine superiore
    bottom=2.5cm, % Margine inferiore (riducilo se vuoi, es. 2cm)
    left=2.5cm,   % Margine sinistro
    right=2.5cm,  % Margine destro
    bindingoffset=0mm, % Se vuoi rilegare, metti es. 5mm
    heightrounded % Arrotonda l'altezza del testo per evitare righe "mozze"
]{geometry}

\raggedbottom






% Formattazione dei titoli di sezione
\setkomafont{chapter}{\Huge\bfseries} % Stile per i capitoli
\RedeclareSectionCommand[beforeskip=1.5em plus .2em minus .1em,afterskip=.8em plus .1em minus .1em]{section}
\RedeclareSectionCommand[beforeskip=1em plus .2em minus .1em,afterskip=.5em plus .1em minus .1em]{subsection}

% -------------------------------------------------------------------
% MATEMATICA
% -------------------------------------------------------------------
\usepackage{amsmath, amssymb, amsthm, mathtools} % Pacchetti base per la matematica
\usepackage{physics}             % Notazione fisica (derivata, bra-ket, etc.)
\usepackage{cancel}              % Per barrare espressioni matematiche
\DeclareSymbolFont{bbold}{U}{bbold}{m}{n}
\DeclareSymbolFontAlphabet{\mathbbold}{bbold}
\newcommand{\id}{\scalebox{1.15}{$\mathbbold{1}$}}

% -------------------------------------------------------------------
% GRAFICA, FIGURE E COLORI
% -------------------------------------------------------------------
\usepackage{graphicx}            % Per includere immagini
\graphicspath{{figures/}}        % Cartella di default per le immagini
\usepackage{subcaption}          % Sotto-didascalie per le figure
\usepackage{tikz}                % Per il disegno vettoriale
\usepackage{pgfplots}            % Per i grafici con TikZ
\pgfplotsset{compat=1.18}         % Imposta la versione di pgfplots
\usepackage[most]{tcolorbox}     % Per creare box colorati
\usepackage{float}

% -------------------------------------------------------------------
% RIFERIMENTI 
% -------------------------------------------------------------------

\usepackage[colorlinks=true, urlcolor=blue, linkcolor=black, citecolor=black]{hyperref} % Link interni ed esterni (deve essere quasi l ultimo)
\usepackage[nameinlink,noabbrev]{cleveref} % Riferimenti "intelligenti"

% -------------------------------------------------------------------
% ELENCHI
% -------------------------------------------------------------------
\usepackage{enumitem}            % Personalizzazione degli elenchi
\usepackage{multicol}            % Per usare pi colonne

% ===================================================================
% AMBIENTI PERSONALIZZATI (Teoremi, Definizioni, etc.)
% ===================================================================

% -------------------------------------------------------------------
% AMBIENTI TEOREMICI
% -------------------------------------------------------------------
\newtheorem{theorem}{Teorema}[section]
\newtheorem{lemma}[theorem]{Lemma}
\theoremstyle{definition}
\newtheorem{definition}{Definizione}[section]
\newtheorem{example}{Esempio}[section]
\newtheorem{exercise}{Esercizio}[section]

\theoremstyle{remark}
\newtheorem*{remark}{Osservazione}
\newtheorem*{solution}{Soluzione}

\newtheoremstyle{attentionstyle}{3pt}{3pt}{\normalfont}{}{\bfseries}{ }{0pt}{#1#3}
\theoremstyle{attentionstyle}
\newtheorem*{attention}{Attenzione!}

\tcbuselibrary{theorems, breakable} % <--- Aggiungi breakable qui

\newtcbtheorem[number within=section]{note}{Nota}{
    enhanced,
    breakable,                      % <--- Questo permette il cambio pagina
    colback=white,
    colframe=black,
    boxrule=0.5pt,
    sharp corners,
    fonttitle=\bfseries,
    description font=\bfseries,
    coltitle=black,
    separator sign none,
    description delimiters={}{},
    attach title to upper,
    after title={\par\medskip},
    left=2mm, right=2mm, top=2mm, bottom=2mm,
}{nt}


% Color box per equazioni importanti

\newtcolorbox{importantbox}{
    colback=blue!5, 
    colframe=blue!50, 
    sharp corners, 
    boxrule=0.5pt,
    % Opzioni per gestire meglio lo spazio intorno alle equazioni
    before skip=10pt,
    after skip=10pt,
    left=2mm, right=2mm, top=2mm, bottom=2mm
}

% ===================================================================
% COMANDI PERSONALIZZATI E IMPOSTAZIONI FINALI
% ===================================================================

% -------------------------------------------------------------------
% COMANDI UTILI
% -------------------------------------------------------------------
\usepackage{marginnote} % Per le note a margine
% --- Gestione Indice delle Lezioni (Versione Forzata) ---
\DeclareNewTOC[
    type=lezione,
    types=lezioni,
    listname={Indice delle Lezioni},
    extension=loz,
]{loz}
\setuptoc{loz}{totoc}
% Questo blocco forza la rimozione degli spazi verticali (addvspace) 
% specificamente per l'indice con estensione .loz
\makeatletter
\BeforeStartingTOC[loz]{%
  \renewcommand*{\addvspace}[1]{}% Annulla il comando che aggiunge spazio tra capitoli
}
\makeatother

\newcommand{\lezione}[2]{%
  \marginnote{\footnotesize\sffamily\textcolor{blue!50!gray}{\textbf{Lez. #1}\\\makebox[\marginparwidth][l]{#2}}}%
  \pdfbookmark[2]{Lezione #1 (#2)}{lez#1}%
  \addcontentsline{loz}{lezione}{\textbf{Lezione #1} (#2)}%
}

% Prodotto interno. Nota: il pacchetto 'physics' offre \expval{} e \bra\ket.
\renewcommand{\ip}[2]{\langle #1, #2 \rangle}

% -------------------------------------------------------------------
% IMPOSTAZIONI FINALI
% -------------------------------------------------------------------
\setcounter{secnumdepth}{3} % Numera fino alle subsubsections
\setcounter{tocdepth}{2}     % Includi fino alle subsections nell indice

% MarkFormat per capitoli (rimuove "Capitolo X:")
\automark{chapter}
\renewcommand*{\chaptermarkformat}{}


% --- Soluzione per i titoli a fine pagina ---
\usepackage{needspace}
\usepackage{etoolbox}

% Prima di ogni \section, controlla se ci sono almeno 6 righe di spazio (~3cm).
% Se non ci sono, manda tutto alla pagina successiva.
\preto\section{\needspace{6\baselineskip}}

% Fa la stessa cosa per le \subsection (richiede un po' meno spazio, es. 4 righe)
\preto\subsection{\needspace{4\baselineskip}}
\preto\subsubsection{\needspace{4\baselineskip}}